\documentclass{article}%
\usepackage[T1]{fontenc}%
\usepackage[utf8]{inputenc}%
\usepackage{lmodern}%
\usepackage{textcomp}%
\usepackage{lastpage}%
\usepackage{geometry}%
\geometry{margin=1in,headheight=14pt,headsep=25pt}%
\usepackage{hyperref}%
\usepackage{enumitem}%
\usepackage{fancyhdr}%
\usepackage{xcolor}%
\usepackage{url}%
\usepackage{breakurl}%
\usepackage{amsmath}%
\usepackage{amssymb}%
\usepackage{mathtools}%
\usepackage{amsthm}%
\usepackage{thmtools}%
%

            \hypersetup{
                colorlinks=true,
                linkcolor=blue,
                filecolor=magenta,
                urlcolor=blue,
                breaklinks=true,
                bookmarks=true,
                pdfborder={0 0 0}
            }
            \urlstyle{same}
            \Urlmuskip=0mu plus 1mu  % URL breaking in nicer spots
            
            % Math configuration
            \everymath{\displaystyle}
            \setlength{\jot}{10pt}  % Increase spacing between equations
            
            \pagestyle{fancy}
            \fancyhf{}
            \rhead{Generated on \today}
            \cfoot{\thepage}
            
            % Better Unicode support
            \usepackage[utf8]{inputenc}
            \usepackage{newunicodechar}
            \usepackage[breaklinks=true]{hyperref}
            \setlength{\itemsep}{1em} % Adds spacing between items
            
            % Configure itemize settings globally
            \setlist[itemize]{nosep}
            \setlist[itemize]{leftmargin=*}
            
            % Remove section numbering
            \renewcommand{\thesection}{}
            \renewcommand{\thesubsection}{}
        %
\lhead{Lecture: 2024{-}10{-}28{-}ApplsLinProg}%
\title{2024{-}10{-}28{-}ApplsLinProg}%
\author{Generated by Scribe.AI}%
\date{\today}%
%
\begin{document}%
\normalsize%
\maketitle%
\section{Slide 1}%
\label{sec:Slide1}%
\subsection{Content}%

\underline{\text{Production Planning}} \\
\text{[M\&D] p.16} \quad 2.3 \text{ Ice Cream All Year Round}
%
\subsection{Description}%
This slide introduces a production planning problem for an ice cream manufacturer. The title "Production Planning" is underlined. A reference to a textbook or other source ([M\&D] p.16) is given, along with the section number (2.3) and title "Ice Cream All Year Round". The main content is a bar chart illustrating monthly ice cream demands ($d_i$, where $i$ represents the month from 1 to 12) in tons. The chart shows the predicted monthly sales of ice cream for the next year, providing data for a linear programming optimization problem. The context of the course (Linear Programming and Optimization Techniques) indicates that this data will be used to formulate a linear program to optimize the ice cream production plan, likely aiming to minimize costs or maximize profits while meeting the predicted demand.%
\subsection{Figures}%
\begin{figure}[h]%
\centering%
\begin{minipage}{\textwidth}%
\centering%
\begin{minipage}{0.15\textwidth}%
\centering%
\includegraphics[width=0.8\textwidth]{/Users/ashoksaravanan/Coding/ScribeLec/server/output/MA 421/lectures/2024-10-28-ApplsLinProg/figures/1.1.png}%
\caption{A bar chart showing monthly demands ($d_i$, $i=1,\dots,12$) for ice cream in tons. The x-axis represents the months (Jan-Dec), and the y-axis represents the sales in tons.}%
\end{minipage}%
\end{minipage}%
\end{figure}

%
\newpage%
\section{Slide 2}%
\label{sec:Slide2}%
\subsection{Content}%

\underline{\text{Production Planning}} \\
\text{[M\&D] p.16} \quad 2.3 \text{ Ice Cream All Year Round}
%
\subsection{Description}%
This slide presents a production planning problem for Icicle World Ltd., an ice cream manufacturer.  The title "Production Planning" is underlined.  A reference to a textbook or other source ([M\&D] p.16) is given, along with the section number (2.3) and title "Ice Cream All Year Round". The main content is a bar chart illustrating the monthly ice cream demand ($d_i$, where $i$ represents the month from 1 to 12) in tons.  The chart shows predicted monthly sales for the next year. A horizontal red line indicates the average monthly demand. The context of the course (Linear Programming and Optimization Techniques) indicates that this data will be used to formulate a linear program to optimize the ice cream production plan, likely aiming to minimize costs or maximize profits while meeting the predicted demand.  The text mentions that the prediction is based on historical data, extensive surveys, and bird observations.%
\subsection{Figures}%
\begin{figure}[h]%
\centering%
\begin{minipage}{\textwidth}%
\centering%
\begin{minipage}{0.15\textwidth}%
\centering%
\includegraphics[width=0.8\textwidth]{/Users/ashoksaravanan/Coding/ScribeLec/server/output/MA 421/lectures/2024-10-28-ApplsLinProg/figures/2.1.png}%
\caption{A bar chart showing monthly demands ($d_i$, $i=1,\dots,12$) for ice cream in tons. The x-axis represents the months (Jan-Dec), and the y-axis represents the sales in tons. A horizontal line indicates the average demand.}%
\end{minipage}%
\end{minipage}%
\end{figure}

%
\newpage%
\section{Slide 3}%
\label{sec:Slide3}%
\subsection{Content}%

\underline{\text{Production Planning}} \\
\text{[M\&D] p.16} \\
(1) \quad x_i - \text{production at month } i \\
(2) \quad s_i - \text{surplus at end of month } i \\
(3) \quad x_i + s_{i-1} = d_i + s_i \\
\min \sum_{i=1}^{12} 50|x_i - x_{i-1}| + \sum_{i=1}^{12} 20 s_i
%
\subsection{Description}%
This slide formulates a linear programming model for a production planning problem. The title "Production Planning" is underlined. A reference to a textbook or other source ([M\&D] p.16) is given.  The model involves: (1) $x_i$ representing the production level in month $i$; (2) $s_i$ representing the surplus inventory at the end of month $i$; and (3) a constraint $x_i + s_{i-1} = d_i + s_i$, where $d_i$ is the demand in month $i$, ensuring that production plus beginning inventory equals demand plus ending inventory. The objective function is to minimize the total cost, which consists of two parts: $\sum_{i=1}^{12} 50|x_i - x_{i-1}|$ representing the cost of changing production levels between consecutive months, and $\sum_{i=1}^{12} 20 s_i$ representing the cost of storing surplus inventory.  The context of the course (Linear Programming and Optimization Techniques) indicates that this model will be solved using linear programming techniques. The absolute value in the cost of changing production levels suggests that this model might require reformulation to be solved using standard linear programming methods.%
\subsection{Figures}%
\begin{figure}[h]%
\centering%
\begin{minipage}{\textwidth}%
\centering%
\begin{minipage}{0.15\textwidth}%
\centering%
\includegraphics[width=0.8\textwidth]{/Users/ashoksaravanan/Coding/ScribeLec/server/output/MA 421/lectures/2024-10-28-ApplsLinProg/figures/3.1.png}%
\caption{No figure present.}%
\end{minipage}%
\end{minipage}%
\end{figure}

%
\newpage%
\end{document}